%!TEX program = xelatex
\title          {Logic in Computer Science - Assignment 7}
\author         {李鹏达 10225101460}
\date{}

\documentclass{article}
\usepackage{amsmath}
\usepackage{../boxproof}
\usepackage{a4wide}
\usepackage{ctex}
\usepackage{enumitem}
\usepackage{amsmath}
\usepackage{amsfonts}
\usepackage{amssymb}
\usepackage{graphicx}
\usepackage{float}
\usepackage{listings}
\usepackage{color}
\usepackage{geometry}
\usepackage{fancyhdr}
\usepackage{titlesec}
\usepackage{tikz}
\usepackage{pgfplots}
\usepackage{pgfplotstable}
\usepackage{booktabs}
\usepackage{caption}
\usepackage{subcaption}
\usepackage{algorithm}
\usepackage{multirow}
\usepackage{array}
\usepackage{enumitem}
\usepackage{hyperref}
\usepackage{fontspec}
\usepackage{algpseudocode}
\usepackage{booktabs}
\usepackage{multirow}
\usepackage{colortbl}
\usepackage{stfloats}
% \usepackage{daymonthyear}

\def\meta#1{\mbox{$\langle\hbox{#1}\rangle$}}
\def\macrowitharg#1#2{{\tt\string#1\bra\meta{#2}\ket}}

{\escapechar-1 \xdef\bra{\string\{}\xdef\ket{\string\}}}

\showboxbreadth 999
\showboxdepth 999
\tracingoutput 1


\let\imp\to

\def\premise{\mathrm{premise}}
\def\assumption{\mathrm{assumption}}
\def\MT{\mathrm{MT\ }}
\def\LEM{\mathrm{LEM}}
\def\intro{\mathrm{i\ }}
\def\elim{\mathrm{e\ }}
\def\introa{\mathrm{i_1\ }}
\def\elima{\mathrm{e_1\ }}
\def\introb{\mathrm{i_2\ }}
\def\elimb{\mathrm{e_2\ }}

\def\lt{<}
\def\eqdef{=}

\def\eps{\mathrel{\epsilon}}
\def\biimplies{\leftrightarrow}
\def\flt#1{\mathrel{{#1}^\flat}}
\def\setof#1{{\left\{{#1}\right\}}}
\let\implies\to
\def\KK{{\mathsf K}}
\let\squashmuskip\relax


\def\pre{\premise}
\def\ass{\assumption}
\def\andea{\land\elima}
\def\andeb{\land\elimb}
\def\andi{\land\intro}
\def\toe{\to\elim}
\def\toi{\to\intro}
\def\ore{\lor\elim}
\def\ore{\lor\elim}
\def\oria{\lor\introa}
\def\orib{\lor\introb}
\def\nege{\neg\elim}
\def\negi{\neg\intro}
\def\bote{\bot\elim}
\def\falle#1{\forall#1\ \elim}
\def\falli#1{\forall#1\ \intro}
\def\exie#1{\exists#1\ \elim}
\def\exii#1{\exists#1\ \intro}

\def\tt{\!-\!}
\allowdisplaybreaks

\lstset{
    language=C,
    basicstyle=\small\ttfamily,
    keywordstyle=\color{blue},
    commentstyle=\color{green!60!black},
    stringstyle=\color{red},
    numbers=left,
    numberstyle=\footnotesize\color{gray},
    stepnumber=1,
    numbersep=5pt,
    backgroundcolor=\color{white},
    showspaces=false,
    showstringspaces=false,
    showtabs=false,
    frame=single,
    tabsize=2,
    breaklines=true,
    breakatwhitespace=false
}

%=======================================================================
\begin{document}
\maketitle
%=======================================================================
Verify the following Hoare triple is partial correct.
\begin{align*}
    &\{n > 0 \land \forall i. \ 0 \leq i < n \to a[i] \in \mathbb{N} \} \\
    &Sort \\
    &\{ \forall i, j. \ 0 \leq i \leq j < n \to a[i] \leq a[j] \}
\end{align*}

The code of Sort is as follows:

~\\
\begin{lstlisting}[language=C,xleftmargin=6em,xrightmargin=22em]
i = 0;
while (i < n-1) {
    j = 0;
    while (j < n-i-1) {
        if (a[j] > a[j+1]) {
            temp = a[j];
            a[j] = a[j+1];
            a[j+1] = temp;
        }
        j = j + 1;
    }
    i = i + 1;
}
\end{lstlisting}

\newpage
\textbf{Solution:}

\begin{align*}
&\text{Denote} \quad \forall k, l. \ k \in [n-i-1..n-2] \times l \in [k+1..n-1] \to a[k] \leq a[l] \\
&\quad\quad\text{as} \quad INV_1; \\
&\quad \forall \ a[j] = \max a[0..j] \quad \text{as} \quad INV_2. \\
\\
& \{ n > 0 \land \forall i. \ 0 \leq i < n \to a[i] \in \mathbb{N} \} & \text{pre} \\
& \{ \forall k, l. \ k \in [n-1..n-2] \land l \in [k+1..n-1] \to a[k] \leq a[l] \}  & \text{implied} \\
1 \ & i = 0; \\
& \{ INV_1 \land i \leq n-1 \} \\
& \{ INV_1 \}  & \text{assignment} \\
2 \ & \text{while } (i < n-1) \ \{ \\
& \{ INV_1 \land i < n-1 \} & \text{inv \& grd} \\
& \{ INV_1 \land i < n \land a[0] = \max[0..0] \}  & \text{implied} \\
3 \ & j = 0; \\
& \{ INV_1 \land i < n-1 \land INV_2 \}  & \text{assignment} \\
4 \ & \quad \text{while } (j < n-i-1) \ \{ \\
& \quad \{ INV_1 \land i < n-1 \land j < n-i-1 \land INV_2 \}  & \text{inv \& grd} \\
5 \ & \quad \quad \text{if } (a[j] > a[j+1]) \ \{ \\
& \quad \quad \{ INV_1 \land i < n-1 \land j < n-i-1 \land INV_2 \land a[j] > a[j+1] \}  & \text{if-statement} \\
6 \ & \quad \quad \quad temp = a[j]; \\
& \quad \quad \{ INV_1 \land i < n-1 \land j < n-i-1 \land INV_2 \land temp > a[j+1] \}  & \text{assignment} \\
& \quad \quad \quad a[j] = a[j+1]; \\
7 \ & \quad \quad \{ INV_1 \land i < n-1 \land j < n-i-1 \land INV_2 \land temp > a[j] \}  & \text{assignment} \\
& \quad \quad \quad a[j+1] = temp; \\
8 \ & \quad \quad \{ INV_1 \land i < n-1 \land j < n-i-1 \land INV_2 \land a[j] < a[j+1] \}  & \text{assignment} \\
9 \ & \quad \quad \} \\
& \quad \quad \{ INV_1 \land i < n-1 \land j < n-i-1 \land INV_2 \land a[j+1] \geq a[j] \}  & \text{if-statement} \\
& \quad \quad \{ INV_1 \land i < n-1 \land j < n-i-1 \land a[j+1] = \max a[0..j+1] \}  & \text{implied} \\
10 \ & \quad \quad j = j + 1; \\
& \quad \quad \{ INV_1 \land i < n-1 \land INV_2 \}  & \text{assignment} \\
11 \ & \quad \} \\
& \quad \{ INV_1 \land i < n-1 \land INV_2 \land j \geq n-i-1 \}  & \text{while} \\
& \quad \{ \forall k, l. \ k \in [n-i-2..n-2] \land l \in [k+1..n-1] \to a[k] \leq a[l] \ \land \\
& i < n-1 \land a[n-i-1] = \max a[0..n-i-1] \}  & \text{implied} \\
12 \ & \quad i = i + 1; \\
& \quad \{ \forall k, l. \ k \in [n-i-1..n-2] \land l \in [k+1..n-1] \to a[k] \leq a[l] \land i < n \}  & \text{assignment} \\
& \quad \{ INV_1 \}  & \text{implied} \\
13 \ & \} \\
& \{ INV_1 \land i \geq n-1 \}  & \text{while} \\
& \{ \forall k, l. \ k \in [0..n-2] \land l \in [k+1..n-1] \to a[k] \leq a[l] \} & \text{implied} \\
& \{ \forall i, j. \ 0 \leq i < j < n \to a[i] \leq a[j] \}  & \text{post} \\
& & \square
\end{align*}
\end{document}

